\documentclass[11pt]{article}
\usepackage[T1]{fontenc}
\usepackage[margin=1in]{geometry}
\usepackage{enumitem}
\usepackage{hyperref}
\setlength{\parindent}{1.5em}
\setlength{\parskip}{6pt}
% =====================================================================
%  EDIT YOUR DETAILS HERE — this ONE file feeds every abstract & deck.
%  (Change a value, recompile; all 36 documents update.)
% =====================================================================
\newcommand{\seminarName}{Aflah Muhammed P}      % <-- your full name (guessed from your email; please verify)
\newcommand{\seminarRoll}{TVE00CS000}            % <-- your university register/roll number
\newcommand{\seminarClassRoll}{00}               % <-- your class roll no (the short one)
\newcommand{\seminarGuide}{Prof.\ <Guide Name>}  % <-- your guide
\newcommand{\seminarDept}{Department of Computer Science and Engineering}
\newcommand{\seminarCollege}{College of Engineering Trivandrum}
\newcommand{\seminarShortCollege}{CET}           % shown in the slide footer, e.g. "Aflah Muhammed P (CET)"
\newcommand{\seminarShortTitle}{B.Tech Seminar}  % middle of the slide footer
\newcommand{\seminarDate}{July 31, 2026}
% ---------------------------------------------------------------------
% Optional: put a logo image named  cet_logo.png  in this folder and it
% will appear on every title slide automatically. If absent, it is skipped.
% =====================================================================

\begin{document}
\begin{center}
{\LARGE\bfseries APIGen-MT: Generating Multi-Turn Tool-Use Data via Agent-Human Simulation}\\[1.1em]
{\large\bfseries Seminar Abstract}\\[0.6em]
{\normalsize \seminarDate}
\end{center}
\vspace{0.6em}
\section*{Abstract}
Large language model (LLM) agents increasingly operate in interactive, multi-turn settings where they must call external tools and APIs, track evolving state, and cooperate with a human user across an entire conversation. Training such agents demands large volumes of high-quality multi-turn tool-use trajectories, yet this data is scarce and difficult to obtain. Human annotation is expensive, slow, and inconsistent, while existing synthetic pipelines, such as single-turn function-calling generators, capture isolated API calls but ignore the dialogue dynamics, partial user intent, and long-horizon dependencies that characterise realistic agentic tasks. Crucially, most approaches lack a rigorous way to verify that generated trajectories are actually correct with respect to ground-truth actions, which limits their reliability as a training signal.

This seminar presents \emph{APIGen-MT}, an agentic two-phase pipeline for generating verified multi-turn tool-use data. In the first phase, a data-generation agent proposes detailed task blueprints with ground-truth actions, which are filtered through format and execution checks and a committee of LLM reviewers, with an iterative feedback loop refining rejected candidates. In the second phase, each validated blueprint is realised as a full conversation through simulated agent-human interplay, retaining only trajectories whose executed actions and outputs match the blueprint. Using 5,000 released trajectories, the authors train the open \emph{xLAM-2} family (1B--70B, built on Llama 3.1/3.2 and Qwen 2.5). xLAM-2-70b reaches 56.2\% on $\tau$-bench and 78.19\% on BFCL v3, surpassing GPT-4o and rivalling Claude 3.5 Sonnet, while smaller models often outperform larger ones in multi-turn settings. The talk covers the pipeline, architecture, results, and limitations.
\section*{Key Reference Papers}
\begin{enumerate}
  \item A. Prabhakar, Z. Liu, M. Zhu, et al., ``APIGen-MT: Agentic Pipeline for Multi-Turn Data Generation via Simulated Agent-Human Interplay,'' arXiv:2504.03601, 2025.
  \item Z. Liu, T. Hoang, J. Zhang, et al., ``APIGen: Automated Pipeline for Generating Verifiable and Diverse Function-Calling Datasets,'' in Proc. NeurIPS, 2024, arXiv:2406.18518.
  \item S. Yao, N. Shinn, P. Razavi, and K. Narasimhan, ``$\tau$-bench: A Benchmark for Tool-Agent-User Interaction in Real-World Domains,'' arXiv:2406.12045, 2024.
\end{enumerate}
\vspace{0.8em}
\noindent\textbf{Submitted By}\\
\seminarName\\
\seminarRoll

\vspace{0.6em}
\noindent\textbf{Guide}\\
\seminarGuide
\end{document}
